\documentclass[11pt,a4paper]{article}
\usepackage[margin=1in]{geometry}
\usepackage{amsmath,booktabs,hyperref,float}

\title{Iteration 021: Mixture of FIR Experts}
\author{Autoresearch Loop}
\date{April 8, 2026}

\begin{document}
\maketitle

\begin{abstract}
A Mixture of Experts model with 4 parallel FIR filters and input-dependent
gating achieves $r=0.375$, not improving over the baseline ($r=0.378$).
The gating network does not learn meaningful expert specialization, and
InstanceNorm at input further reduces performance.
\end{abstract}

\section{Hypothesis}
Different test subjects need different spatial filters. An MoE model with
$K=4$ experts learns multiple FIR filters and selects them based on input
channel statistics, adapting to different head geometries without subject labels.

\section{Method}
Architecture: InstanceNorm $\to$ 4 parallel FIR experts (each $27 \to 12$,
7 taps) $\to$ gating network (channel means $\to$ softmax weights) $\to$
weighted combination. All experts CF-initialized with small perturbation.
Combined loss + corr validation + 15\% channel dropout.

\section{Results}
\begin{table}[H]
\centering
\begin{tabular}{lrrr}
\toprule
Model & Mean $r$ & Std $r$ & SNR (dB) \\
\midrule
Combined loss (iter017) & 0.378 & 0.078 & 0.61 \\
MoE (4 experts) & 0.375 & 0.077 & 0.63 \\
\bottomrule
\end{tabular}
\end{table}

\section{Analysis}
Two compounding issues: (1) InstanceNorm removes useful amplitude information
(consistent with iter020), and (2) the experts start from nearly identical
CF initializations, making it difficult for the gating network to learn
meaningful specialization. With only 12 training subjects, there may not be
enough diversity to drive expert divergence.

The slightly higher SNR but lower $r$ matches the InstanceNorm pattern seen
across iterations 020--025: better amplitude matching but worse waveform shape.

\section{Next}
Test augmentation techniques (mixup, noise, TTA) without InstanceNorm.

\end{document}
